% ══════════════════════════════════════════════════════════════
%  SECTION 6: HOW THE GAME CHANGED — ON-PITCH EFFECTS
%  Draft v3 — all numbers from results_v3_onpitch.json
% ══════════════════════════════════════════════════════════════
\section{How the Game Changed: On-Pitch Effects}
\label{sec:onpitch}

With the first stage established (Section~\ref{sec:firststage}) and the compliance mechanism documented (Section~\ref{sec:compliance}), we turn to the central question: what happened on the pitch?\footnote{Our three pre-specified primary outcomes are second-half stoppage time (first stage), total goals per match, and goals after the 90th minute. All other on-pitch outcomes are exploratory. We report unadjusted $p$-values throughout; readers should apply multiple-comparison caution to exploratory results.} Total goals were unchanged. The per-minute intensity of play held steady. The directive's signature effect was not the creation of new competitive action but its \emph{redistribution} toward the end of matches---more late drama, the same overall output.


\subsection{Goals: Displacement, Not Creation}
\label{subsec:goals}

We begin with the outcome most closely watched by fans, managers, and broadcasters. Table~\ref{tab:onpitch_effects_v3} reports the difference-in-differences estimate for total goals per match. The point estimate is $\hat{\beta} = -0.033$ (SE $= 0.038$, $p = 0.39$, $N = 45{,}433$): a precise zero. Pre-directive, matches averaged 2.649 total goals; post-directive, 2.775. The raw increase of 0.126 goals is fully explained by secular trends common to treated and control leagues, leaving no residual treatment effect.

This null is well-powered. With SE $= 0.038$, the design can detect an increase
of 0.11 goals per match (4\% of the pre-treatment mean) at 80\% power. Moreover,
we can reject not only the null of no effect but also the mechanical prediction
that more minutes at the historical scoring rate should produce approximately
$+0.044$ additional goals ($t = (-0.033 - 0.044)/0.038 = -2.03$, $p = 0.042$).
The offset is near-complete, and the data are inconsistent with a simple ``more
time, more goals'' model.

The null on total goals coexists with a highly significant increase in \emph{late} goals. In a pre-post comparison within treated leagues, goals scored after the 90th minute rose by 0.030 per match ($\text{SE} = 0.007$, $p < 0.001$), from a pre-directive mean of 0.153 to a post-directive mean of 0.183---a 20\% relative increase. The per-match rate in the 90+ window increased by 35.5\%, dwarfing the modest changes in earlier time bins. Meanwhile, goals in the 76--84 minute window \emph{fell} by 8.6\%, and the 85--89 window declined by 2.7\%.

This pattern---gains at the very end, losses in the penultimate windows---is the hallmark of displacement. Any football supporter recognizes the dynamic: when a team trailing by one goal knows that seven or eight minutes of stoppage remain rather than three, the desperation that previously peaked at minute 82 can be deferred. Managers hold their final substitution longer; central defenders push into the attacking half later. The directive did not create additional goals; it shifted them later. Scoring opportunities that would have materialized in the waning minutes of regular time were redistributed into stoppage time. The match as a productive unit generated the same output; only the timing changed.

The result-distribution DiD estimates reinforce this conclusion. The probability of a draw ($\hat{\beta} = 0.003$, $p = 0.80$), home win ($\hat{\beta} = -0.008$, $p = 0.49$), and away win ($\hat{\beta} = 0.005$, $p = 0.61$) are all precisely estimated nulls. The directive did not make home advantage more or less valuable; it did not generate more draws or fewer. The competitive outcome of the match was invariant to the treatment.

\begin{figure}[H]
\centering
\includegraphics[width=\textwidth]{fig_goal_timing_v3}
\caption{Goals per match by time bin, Big Five leagues, pre- vs.\ post-directive. The post-directive period shows more goals after the 90th minute and fewer in the 76--89 window, consistent with displacement rather than creation.}
\label{fig:goaltiming}
\end{figure}


\subsection{The Quantity--Intensity Decomposition}
\label{subsec:decomposition}

To formalize the displacement result, we decompose the change in each match-level outcome into two channels. Let $Y$ denote an outcome (goals, fouls, cards) observed over the match. We write $Y = r \times T$, where $r$ is the per-minute rate and $T$ is effective playing time. The observed change between regimes decomposes as:
\begin{equation}
\label{eq:decomposition}
    \Delta Y = \underbrace{r_{\text{pre}} \cdot \Delta T}_{\text{quantity effect}} + \underbrace{\Delta r \cdot T_{\text{post}}}_{\text{intensity effect}}
\end{equation}
The \emph{quantity effect} measures how much of the change in $Y$ is attributable to more playing time at the pre-treatment per-minute rate; the \emph{intensity effect} captures any shift in that rate.\footnote{The exact identity is $\Delta Y = r_{\text{pre}} \cdot \Delta T + \Delta r \cdot T_{\text{pre}} + \Delta r \cdot \Delta T$. Our formulation assigns the interaction term to the intensity effect through post-period weighting. Qualitative conclusions are unchanged under pre-period weighting.}

\paragraph{Goals after 90 minutes.} The per-minute scoring rate in second-half stoppage time was 0.034 before the directive and 0.046 after---a 35.5\% increase in the raw rate. This apparent intensification is misleading: it reflects the composition of longer stoppages rather than a change in the fundamental rate of play during any given stoppage minute. When stoppages extend from four minutes to six, the additional two minutes include the chaotic, open-play phase in which trailing teams have committed their centre-backs forward and defensive shape has dissolved. These minutes are mechanically high-scoring, but they represent a selection effect, not a behavioral change.

The raw pre-post increase in treated leagues is 0.060 goals per match in the post-90th-minute window. Applying the decomposition in Equation~\ref{eq:decomposition}, 87\% of this increase is attributable to the quantity effect (more minutes at the pre-treatment rate) and 13\% to an intensity residual, with the intensity component not significantly different from zero. This is the purest form of displacement: the directive purchased additional minutes, and those minutes produced goals at exactly the historical rate. No new attacking ingenuity was required; no tactical revolution occurred. Teams simply had more time to do what they were already doing.

\paragraph{Total goals.} The per-match-minute scoring rate was 0.028 goals per minute before and after ($p = 0.47$). The observed raw increase within treated leagues ($+0.126$ per match) is modest relative to the pre-directive mean of 2.649 goals per match.% TODO: recompute decomposition shares with correct base of 0.126
After accounting for control-group trends via DiD, the raw increase disappears entirely ($\hat{\beta} = -0.033$, $p = 0.39$), confirming that the treated-league trend tracks the secular rise in scoring across all leagues rather than reflecting a directive-specific effect.

\paragraph{Total fouls.} The foul decomposition reveals a behavioral shift. The per-minute foul rate declined significantly from 0.254 to 0.237 ($p < 0.001$). The mechanical quantity effect of 1.6 additional minutes predicts $+0.40$ fouls; instead, fouls fell. However, the DiD coefficient for total fouls is \emph{positive} ($+1.244$, $p < 0.001$; Table~\ref{tab:onpitch_effects_v3}), indicating that control leagues experienced an even steeper secular decline. The intensity drop in treated leagues therefore reflects a league-wide trend rather than a directive-specific behavioral response.

\paragraph{Yellow cards.} Yellow cards show a near-zero net change ($+0.014$
per match) masking offsetting components: a mechanical increase of $+0.065$
expected from more playing time, offset by a behavioral residual of $-0.051$.
The per-minute yellow card rate declined modestly (0.042 to 0.041, $p = 0.009$).
This is consistent with the secular decline in cautions visible across European
football---a trend often attributed to improvements in refereeing consistency
and the professionalization of player conduct---rather than a treatment-specific
response.

\paragraph{Summary.} Table~\ref{tab:decomposition_v3} collects all
decomposition results. Two findings stand out. First, the per-minute scoring
rate---both in stoppage time and across the full match---was unchanged,
confirming displacement as the dominant mechanism for goals. Second, the
per-minute foul rate dropped significantly within treated leagues, but the
DiD attributes this to secular trends rather than the directive
($\hat{\beta}_{\text{fouls}} = +1.244$, $p < 0.001$). Figure~\ref{fig:decomposition_v3}
in Appendix~\ref{app:figures} provides a visual decomposition.

\begin{table}[htbp]
\centering
\caption{Quantity--Intensity Decomposition}
\label{tab:decomposition_v3}
\begin{threeparttable}
\begin{tabular}{lcccccc}
\toprule
Outcome & $\Delta$ Stoppage & Observed $\Delta$ & Quantity & Intensity & Qty Share & Rate $p$ \\
\midrule
  Total goals & +1.56 & +0.065 & +0.044 & +0.020 & 69\% & 0.468 \\
  Total fouls & +1.56 & -0.990 & +0.396 & -1.387 & $-$40\% & 0.000*** \\
  Yellow cards & +1.56 & +0.014 & +0.065 & -0.051 & 461\% & 0.009** \\
  Goals after 90' & +1.56 & +0.053 & +0.046 & +0.007 & 87\% & 0.685 \\
\bottomrule
\end{tabular}
\begin{tablenotes}[flushleft]\footnotesize
\item \textit{Notes:} Quantity channel = $\Delta$stoppage $\times$ pre-treatment per-minute rate. Intensity = observed $-$ expected. Bootstrap 95\% CIs (1,000 iterations). Rate $p$-value from Welch $t$-test of per-minute rates pre vs.\ post. For full-match outcomes (total goals, fouls, cards), the per-minute rate uses total match minutes ($90 + \text{stoppage}_{1H} + \text{stoppage}_{2H}$) as the denominator. For goals after 90', the rate uses second-half stoppage minutes only. Negative quantity shares arise when the quantity and observed changes have opposite signs.
\end{tablenotes}
\end{threeparttable}
\end{table}



\subsection{Game-Changing Goals: When Results Hang in the Balance}
\label{subsec:game_changing}

For football fans, not all goals are created equal. A goal in the 12th minute adjusts expectations; a goal in the 93rd minute that turns a defeat into a draw sends 40,000 people into delirium. The directive's displacement of goals toward the end of matches increases the frequency of precisely these high-salience moments---goals scored late that alter the match result.

We define a \emph{game-changing goal} as any goal scored after the 85th minute that changes which team earns points: an equalizer that converts a loss into a draw, a go-ahead goal that converts a draw into a win, or a second-equalizer that negates an opponent's earlier lead. In football culture, these are the moments around which seasons pivot. A late equalizer against a title rival can be worth as much as three points in the final standings when it also destroys the opponent's momentum. A 94th-minute winner in a relegation six-pointer enters the folklore of a club. These goals carry weight disproportionate to their statistical frequency.

The pre-post estimate for game-changing goals after minute 85 is $\hat{\beta} = 0.035$ ($\text{SE} = 0.007$, $p < 0.001$, $N = 14{,}441$).\footnote{Because minute-level data are available only for Big Five leagues and the Championship, this estimate uses a pre-post comparison within treated leagues with league and season fixed effects, not the full DiD.} Pre-directive, 17.9\% of matches featured at least one such goal; post-directive, 21.4\%---a relative increase of approximately 20\%. One additional match in every six that previously ended quietly now featured a late result-altering goal.

Decomposing the effect, late equalizers rose by 1.2 percentage points ($p = 0.006$), from a base rate of 5.3\% to 6.5\%. Late winners---goals that flip a draw into a victory or overturn a deficit---increased by 1.3 percentage points ($p = 0.013$), from 7.5\% to 8.8\%. Both channels contributed to the drama dividend. The split is roughly symmetric: the directive did not disproportionately favor trailing teams or leaders; it simply created more late decisive moments of all types.

Narrowing to the 90+ window alone, goals that changed the result rose by 3.0 percentage points ($p < 0.001$), from 7.3\% to 10.3\%---a 42\% relative increase in the probability that stoppage time alone reshapes the final score. In concrete terms, one additional match in every ten now sees its result altered during stoppage time.

This is the directive's most palpable on-pitch legacy. A supporter leaving the stadium at minute 88 in the post-directive era forgoes a meaningfully higher probability of witnessing the decisive moment. For broadcasters, these late swings are among the highest-value events in football: they sustain audience retention through the final minutes and generate the highlights that drive post-match engagement. The directive did not make football more productive; it made football more \emph{dramatic} (Figure~\ref{fig:game_changing}).

\begin{figure}[H]
\centering
\includegraphics[width=\textwidth]{fig_game_changing}
\caption{Probability that a goal after the 85th minute changes the match result, pre- vs.\ post-directive. The post-directive increase reflects both late equalizers and late winners.}
\label{fig:game_changing}
\end{figure}


\subsection{Goal Timing Profile}
\label{subsec:goal_timing}

Figure~\ref{fig:goaltiming} reports per-minute scoring rates across nine time bins, providing a complete temporal map of the directive's effects. The first-half bins (1--15, 16--30, 31--45) show virtually no change: per-minute rates of 0.027--0.029 before and after, with percentage changes of $-0.5\%$, $-1.7\%$, and $-2.5\%$, respectively. This stability is expected: the directive primarily targeted second-half stoppage time, and the first 45 minutes of regulation play were structurally unaffected. First-half stoppage time (45+) shows a 14.3\% increase in per-minute rate (0.068 to 0.078), likely reflecting parallel enforcement of stoppage-time norms in the first half, though this period was not the directive's primary focus.

The early second half (46--60, 61--75) is similarly stable. The 46--60 window shows a per-minute rate of 0.025 pre and 0.026 post ($+3.1\%$); the 61--75 window holds at 0.033 ($-0.5\%$). These are the minutes in which most second-half goals are scored (the 61--75 bin produces the highest per-match total of any window, at 0.495 goals), and they were entirely unaffected by the directive.

The action concentrates in the final 15 minutes plus stoppage. The 76--84 window shows a notable \emph{decline} in per-minute scoring ($-8.6\%$, from 0.034 to 0.031)---the largest negative shift in any bin. This is consistent with teams adjusting their pressing behavior in anticipation of longer stoppages. In the pre-directive regime, a team trailing 1--0 at minute 78 faced perhaps four minutes of stoppage time; in the post-directive regime, it faces six or seven. The reduced urgency at minute 80 manifests as a lower per-minute scoring rate in the 76--84 window---not because the quality of play deteriorated, but because teams rationally deferred their final push. The 85--89 window declines modestly ($-2.7\%$), suggesting the behavioral adjustment is concentrated in the late-seventies and early-eighties.

The 90+ window then surges: per-minute scoring rises from 0.034 to 0.046, a 35.5\% increase---the largest movement in any time bin by an order of magnitude. This is not, however, evidence of heightened intensity in stoppage time per se. The per-minute rate in 90+ minutes is comparable to rates in regular second-half play (0.026--0.033). The elevated rate partly reflects compositional effects: the additional stoppage minutes disproportionately feature open, stretched play as trailing teams commit numbers forward and defensive structures fatigue. Centre-backs who have played 93 minutes close gaps less quickly; full-backs who have pushed forward for a late equalizer leave space behind them.

The overall pattern confirms that stoppage time is not inherently more or less productive than regular time---it is simply more of the same football, played at the same fundamental rate. This has a practical implication for league scheduling and broadcasting: the additional minutes produce match content indistinguishable in quality from the minutes they displaced, not a qualitatively different product.


\subsection{Fouls and Cards}
\label{subsec:fouls_cards}

Did the directive alter the strategic calculus of fouling? In the pre-directive
regime, the ``tactical foul'' was one of football's most reliable
time-management tools: a defender pulls back an attacker at the halfway line in
the 87th minute, concedes a free kick, and burns 30 seconds while the referee
administers a yellow card. If the directive made referees more diligent about
adding back time lost to such incidents, the expected return to tactical fouling
should decline, and rational teams should foul less.

The DiD estimates tell a more nuanced story than this simple prediction. Total
fouls show a \emph{positive} treatment effect ($\hat{\beta} = 1.244$,
SE $= 0.130$, $p < 0.001$, $N = 39{,}130$), meaning treated leagues experienced
a \emph{smaller} decline in fouls than control leagues. The pre-directive mean
was 25.8 fouls per match in treated leagues; both groups saw secular declines,
but the treated-league decline was 1.2 fouls per match shallower. Home fouls
($\hat{\beta} = 0.596$, $p < 0.001$) and away fouls ($\hat{\beta} = 0.648$,
$p < 0.001$) contributed roughly equally, suggesting no asymmetric effect of
the directive on home-team or away-team fouling strategies.

Yellow cards, by contrast, show no significant treatment effect. Home yellows
($\hat{\beta} = -0.053$, SE $= 0.031$, $p = 0.085$) and away yellows
($\hat{\beta} = -0.020$, SE $= 0.032$, $p = 0.54$) are both near zero. Total
cards are insignificant ($\hat{\beta} = -0.053$, $p = 0.30$; pre-directive mean
$= 4.27$). Red cards are similarly null (home: $\hat{\beta} = 0.009$,
$p = 0.19$; away: $\hat{\beta} = 0.011$, $p = 0.15$). The pre-post estimates
from event-level FBRef data corroborate: total fouls declined by 0.93 per match
($p < 0.001$) within treated leagues, while total cards were unchanged
($\hat{\beta} = 0.026$, $p = 0.52$).

The contrast between fouls and cards is instructive. Referees did not become more or less punitive; the card rate per foul was stable. The relative increase in fouls in treated leagues may reflect the mechanical effect of longer matches---more playing time generates more fouls at a constant per-minute rate---or a subtle behavioral response in which teams, knowing that wasted time will be added back, feel less need to economize on fouling. The ``professional foul'' at the halfway line to prevent a counter-attack carries a lower effective cost when the time it consumes will be restored.

In either case, the directive did not trigger the kind of disciplinary escalation that critics feared when the policy was announced. More time did not mean more violence, more red cards, or more cynical play. The disciplinary ecology of the match---the balance of tactical fouls, cautions, and dismissals that structures the flow of professional football---was preserved.


\subsection{Competitive Balance}
\label{subsec:competitive_balance}

A natural concern is that displacing goals toward the end of matches might systematically advantage certain types of teams. Clubs with deeper squads, superior fitness, or more aggressive late-game tactics could disproportionately benefit from longer stoppages. If so, the directive might widen the gap between rich and poor clubs---a material concern for league structures built on the promise of competitive uncertainty.

We test for changes in competitive balance using two standard measures: the standard deviation of final league points, which captures overall dispersion, and the normalized Herfindahl--Hirschman Index (NHHI) of points concentration, which is sensitive to dominance by a small number of clubs. Neither measure shows a statistically significant change ($\hat{\beta}_{\text{SD}} = 0.172$, SE $= 0.685$, $p = 0.80$; $\hat{\beta}_{\text{NHHI}} = 0.0007$, SE $= 0.0005$, $p = 0.14$; $N = 138$ league-seasons). Pre-directive, the standard deviation of league points was 15.5; post-directive, 15.5. The NHHI was 0.0056 before and 0.0063 after, with the small increase fully absorbed by league fixed effects. The league table's shape was preserved.

The test has 80\% power to detect a change of 1.9 points in the standard deviation (12.4\% of the pre-directive mean), sufficient to rule out large disruptions but not subtle rebalancing. This precision is adequate for the concern at hand: the worry among league officials was not a marginal tilt but a structural shift that would erode competitive uncertainty---the property that makes league football commercially viable.

The null is consistent with the displacement interpretation. If the directive redistributed goals \emph{across time within matches} without changing total goals, and if the redistribution was not systematically correlated with team quality, then league-level competitive balance should be unaffected. The additional late goals reshuffled individual match results---creating winners and losers in specific fixtures---without tilting the season-long distribution of points toward stronger or weaker clubs.

This result should reassure league administrators. The fear that longer stoppages would entrench the dominance of super-clubs---whose deeper squads and superior conditioning might allow them to exploit extended playing time---finds no support in the data. Over a 38-match season, the law of large numbers appears to neutralize any match-level asymmetry: a late equalizer conceded by a top club in one fixture is offset by a late winner snatched in another. The competitive structure of European football survived the directive intact (Figure~\ref{fig:competitive_balance_v3}).

\begin{figure}[H]
\centering
\includegraphics[width=\textwidth]{fig_competitive_balance_v3}
\caption{Competitive balance measures (standard deviation of points and NHHI), treated vs.\ control leagues, pre- and post-directive. Neither measure shows a significant treatment effect.}
\label{fig:competitive_balance_v3}
\end{figure}


\subsection{Distributional Effects}
\label{subsec:distributional}

We examine the full distribution of stoppage time using quantile treatment effects (QTEs). Figure~\ref{fig:quantile_effects_v3} displays the estimated treatment effect at the 10th, 25th, 50th, 75th, and 90th percentiles. The effect is strikingly uniform: approximately 1.0 minute at every quantile ($\hat{\beta}_{10} = 1.00$, $\hat{\beta}_{25} = 1.00$, $\hat{\beta}_{50} = 1.00$, $\hat{\beta}_{75} = 1.00$, $\hat{\beta}_{90} = 1.00$; all $p < 0.001$). The \emph{entire} distribution shifted rightward by approximately one minute, preserving its shape. This is, in distributional terms, a pure location shift.

\begin{figure}[H]
\centering
\includegraphics[width=\textwidth]{fig_quantile_effects_v3}
\caption{Quantile treatment effects on second-half stoppage time. Each point shows the estimated DiD effect at a given percentile. The approximately uniform effect of 1.0 minute across all quantiles indicates a pure location shift rather than a change in distributional shape.}
\label{fig:quantile_effects_v3}
\end{figure}

This uniformity rules out several alternative models. It is inconsistent with a ``floor effect'' in which the directive primarily eliminated very short stoppages while leaving longer ones unchanged. It is inconsistent with a ``drama effect'' in which referees selectively extended close matches but not blowouts. And it is inconsistent with a ``compliance ceiling'' in which only the most egregiously short pre-directive stoppages were corrected. Instead, the data support the simplest possible model: referees added approximately one minute to every match regardless of where that match fell in the pre-directive distribution. The directive was implemented as a uniform bureaucratic adjustment, not a context-dependent behavioral response.

Despite the uniform shift in quantiles, the directive had asymmetric effects on the probability of extreme stoppages---a mechanical consequence of the nonlinear relationship between a location shift and tail probabilities. When the entire distribution moves right by one minute, the mass that crosses any fixed threshold depends on the density at that threshold. Near the center, where density is high, a one-minute shift moves a large share of matches across the threshold; in the thin tails, few additional matches cross.

The probability that stoppage time exceeded 7 minutes rose from 14.8\% to 39.2\% in treated leagues (DiD: $\hat{\beta} = 0.126$, $p < 0.001$), transforming what had been a roughly one-in-seven event into a two-in-five event. The probability of exceeding 10 minutes rose from 1.8\% to 7.8\% (DiD: $\hat{\beta} = 0.029$, $p < 0.001$), a fourfold increase from a low base. For fans and stadium operators accustomed to knowing when a match would end---and for the transport planners, policing coordinators, and broadcasters whose schedules depend on predictable match durations---the post-directive era introduced genuine uncertainty about the length of play.


\subsection{Summary of On-Pitch Effects}

\begin{table}[htbp]
\centering
\caption{On-Pitch Effects of the Stoppage-Time Directive}
\label{tab:onpitch_effects_v3}
\small
\begin{threeparttable}
\begin{tabular}{lccccc}
\toprule
 & $\hat{\beta}$ & SE & Pre-mean & Post-mean & N \\
\midrule
\multicolumn{6}{l}{\textit{Panel A: Goal Effects}} \\
  Goals in 90+ min & +0.0302*** & (0.0069) & 0.153 & 0.183 & 16,212 \\
  Result changed in stoppage & +0.0303*** & (0.0053) & 0.073 & 0.103 & 14,441 \\
  Late equalizer & +0.0120** & (0.0044) & 0.053 & 0.065 & 14,441 \\
  Late winner & +0.0126* & (0.0050) & 0.075 & 0.088 & 14,441 \\
  Total goals per match & -0.0328 & (0.0379) & 2.649 & 2.775 & 45,433 \\
\midrule
\multicolumn{6}{l}{\textit{Panel B: Result Distribution}} \\
  Draw rate & +0.0026 & (0.0100) & 0.264 & 0.261 & 45,433 \\
  Home win rate & -0.0079 & (0.0114) & 0.436 & 0.432 & 45,433 \\
  Away win rate & +0.0054 & (0.0106) & 0.300 & 0.307 & 45,433 \\
\midrule
\multicolumn{6}{l}{\textit{Panel C: Fouls and Cards}} \\
  Total fouls & +1.2441*** & (0.1298) & 25.757 & 24.020 & 39,130 \\
  Total cards & -0.0531 & (0.0516) & 4.269 & 4.307 & 40,469 \\
  Home yellow cards & -0.0530 & (0.0308) & 1.904 & 1.912 & 40,469 \\
  Away yellow cards & -0.0199 & (0.0324) & 2.157 & 2.203 & 40,469 \\
  Home red cards & +0.0091 & (0.0069) & 0.092 & 0.089 & 40,469 \\
  Away red cards & +0.0106 & (0.0074) & 0.116 & 0.104 & 40,469 \\
\midrule
\multicolumn{6}{l}{\textit{Panel D: Competitive Balance}} \\
  SD of points & +0.1717 & (0.6848) & 15.476 & 15.469 & 138 \\
  Normalized HHI & +0.0007 & (0.0005) & 0.006 & 0.006 & 138 \\
\bottomrule
\end{tabular}
\begin{tablenotes}
\small
\item Notes: OLS regressions with league and season FE. HC1 robust SEs in parentheses (the first-stage specification in Table~\ref{tab:first_stage} uses league-clustered SEs; qualitative conclusions are unchanged).
$^{***}p<0.001$, $^{**}p<0.01$, $^{*}p<0.05$.
Total goals and Panels B--D use match-level DiD (treated vs control leagues).
Goals in 90+ min, result changed in stoppage, late equalizer, and late winner use pre-post comparisons within treated leagues with minute-level data.
Panel D uses league-season level DiD.
\end{tablenotes}
\end{threeparttable}
\end{table}

Table~\ref{tab:onpitch_effects_v3} collects the principal estimates. The directive changed the temporal distribution of goals without changing their total quantity (DiD: $-0.033$, $p = 0.39$), raised the frequency of late game-changing moments by 20\% ($+0.035$, $p < 0.001$), did not reduce fouls relative to control leagues (DiD: $+1.244$, $p < 0.001$), and left match outcomes---draws, home wins, away wins---statistically unchanged. Competitive balance, measured by the standard deviation of league points and the NHHI, was unaffected.

The game got longer, late drama became more common, and the competitive
equilibrium was preserved. From an economist's perspective, the directive
functioned as a redistribution of existing match output across time---a temporal
reallocation that increased the entertainment value of the final minutes without
altering the match's aggregate production. The quantity-intensity decomposition
confirms that 87\% of the increase in late goals is attributable to more minutes,
not more intense play. The distributional analysis shows a pure location shift
of approximately one minute at every quantile, ruling out selective enforcement
or context-dependent referee behavior.

From a football perspective, the story is simpler: the same football, played for
longer, with the decisive moments arriving later. The supporter's experience
changed---more uncertainty about when the final whistle would blow, more
late-match tension, more reasons to stay until the end---but the fundamental
character of the competition was preserved. We turn next to whether financial
markets priced this change.
